\documentclass[12pt,a4paper]{article}
\usepackage[utf8]{inputenc}
\usepackage{amsmath,amssymb}
\usepackage{booktabs}
\usepackage[margin=2.5cm]{geometry}

\begin{document}



\textbf{Ya}, $Z = X + Y$ mempunyai fpm. Karena $Z$ hanya mengambil nilai berhingga, yaitu $0$, $1$, dan $2$, maka
\[
M_Z(t) = E\!\left[e^{tZ}\right]
\]
merupakan jumlahan berhingga sehingga selalu ada untuk setiap $t \in \mathbb{R}$.

\subsection*{Langkah 1: Distribusi $Z$}
\begin{center}
\begin{tabular}{ccc}
\toprule
$(x,y)$ & $z = x+y$ & $p(x,y)$ \\
\midrule
$(0,0)$ & 0 & 0{,}4 \\
$(0,1)$ & 1 & 0{,}1 \\
$(1,0)$ & 1 & 0{,}2 \\
$(1,1)$ & 2 & 0{,}3 \\
\bottomrule
\end{tabular}
\end{center}
Sehingga
\begin{align*}
P(Z=0) &= 0{,}4, \\
P(Z=1) &= 0{,}1 + 0{,}2 = 0{,}3, \\
P(Z=2) &= 0{,}3,
\end{align*}
dengan $0{,}4 + 0{,}3 + 0{,}3 = 1$.

\subsection*{Langkah 2: Menghitung fpm}
\begin{align*}
M_Z(t) &= E\!\left[e^{tZ}\right] = \sum_{z} e^{tz}\, P(Z=z) \\
&= e^{0}(0{,}4) + e^{t}(0{,}3) + e^{2t}(0{,}3),
\end{align*}
sehingga
\[
\boxed{M_Z(t) = 0{,}4 + 0{,}3\,e^{t} + 0{,}3\,e^{2t}, \qquad t \in \mathbb{R}.}
\]
Secara setara, fpm dapat dihitung langsung dari distribusi bersama:
\[
M_Z(t) = E\!\left[e^{t(X+Y)}\right] = \sum_{x}\sum_{y} e^{t(x+y)}\, p(x,y),
\]
dan hasilnya sama.

\subsection*{Catatan: $X$ dan $Y$ tidak independen}
Rumus $M_{X+Y}(t) = M_X(t)\,M_Y(t)$ hanya berlaku jika $X$ dan $Y$ independen. Di sini
\[
P(X=0) = 0{,}5, \qquad P(Y=0) = 0{,}6,
\]
sehingga
\[
P(X=0)\,P(Y=0) = 0{,}30 \neq 0{,}4 = P(X=0,\,Y=0).
\]
Jadi $X$ dan $Y$ tidak independen. Jika rumus perkalian tetap dipakai, diperoleh
\[
(0{,}5 + 0{,}5e^{t})(0{,}6 + 0{,}4e^{t}) = 0{,}3 + 0{,}5e^{t} + 0{,}2e^{2t},
\]
yang berbeda dari fpm yang benar.

\subsection*{Pengecekan}
\[
M_Z'(t) = 0{,}3e^{t} + 0{,}6e^{2t}
\quad\Longrightarrow\quad
M_Z'(0) = 0{,}3 + 0{,}6 = 0{,}9.
\]
Hasil ini sama dengan $E[X] + E[Y] = 0{,}5 + 0{,}4 = 0{,}9$, sehingga fpm yang diperoleh konsisten.

\end{document}
